\documentclass[conference]{IEEEtran}
\IEEEoverridecommandlockouts
\usepackage{cite}
\usepackage{amsmath,amssymb,amsfonts}
\usepackage{algorithmic}
\usepackage{graphicx}
\usepackage{textcomp}
\usepackage{xcolor}
\usepackage{bbm}
\usepackage{booktabs}
\usepackage{multirow}
\usepackage{url}
\usepackage{tabularx}
\def\BibTeX{{\rm B\kern-.05em{\sc i\kern-.025em b}\kern-.08em
    T\kern-.1667em\lower.7ex\hbox{E}\kern-.125emX}}
\begin{document}

\title{SM-SIP: Semantic and Multilingual Salient Information Prompting for Abstractive Summarization}

\author{\IEEEauthorblockN{Marc'Antonio Lopez}
\IEEEauthorblockA{\textit{Politecnico di Torino} \\
Torino, Italy \\
marcantonio.lopez@studenti.polito.it}
\and
\IEEEauthorblockN{Chiara Lunazzi}
\IEEEauthorblockA{\textit{Politecnico di Torino} \\
Torino, Italy \\
chiara.lunazzi@studenti.polito.it}
\and
\IEEEauthorblockN{Luca Desderi}
\IEEEauthorblockA{\textit{Politecnico di Torino} \\
Torino, Italy \\
luca.desderi@studenti.polito.it}
}

\maketitle

\begin{abstract}
We present SM-SIP (Semantic \& Multilingual Salient Information Prompting), an extension of SigExt for controllable abstractive summarization. While LLMs generate fluent text, they may omit key facts or hallucinate unsupported content. SIP mitigates this by injecting extracted salient information into the prompt, but its original supervision relies on lexical fuzzy matching. We replace this step with \textbf{semantic supervision} based on Sentence-BERT cosine similarity, and we add a \textbf{multilingual extension} to Italian (WITS). The extractor is trained with Longformer-based backbones and paired with Llama-3.1-8B-Instruct for generation. On Italian, SM-SIP is stable across configurations (average BERTScore \(\approx 0.66\)) and reaches up to 49\% KIR, with the 10k setting performing best, indicating strong data efficiency. These results support semantic label generation as a practical improvement over lexical matching and show that SIP-style controllability transfers effectively to a morphologically rich language. Code and resources are available at \url{https://github.com/LucaDesderi/DNLPProj.git}.
\end{abstract}

\begin{IEEEkeywords}
abstractive summarization, salient information prompting, semantic supervision, multilingual NLP, large language models, Italian text summarization
\end{IEEEkeywords}

%% ============================================================
%% SECTION 1: INTRODUCTION
%% ============================================================
\section{Introduction}
Large Language Models (LLMs) have significantly improved abstractive summarization quality, but controllability remains a core issue: generated summaries can miss critical facts or introduce unsupported statements~\cite{maynez2020faithfulness, cao2022hallucination}. Salient Information Prompting (SIP)~\cite{xu2024sip} addresses this by separating extraction and generation: a dedicated extractor (SigExt) identifies salient evidence that is injected into the generation prompt as soft constraints.

A key limitation of the original pipeline is supervision quality. SigExt labels are derived via character-level fuzzy matching, which is brittle under paraphrase and morpho-syntactic variation. This can propagate noisy labels to training and reduce semantic generalization.

We propose \textbf{SM-SIP}, with two extensions:

\begin{enumerate}
    \item \textbf{Extension 1: Semantic Supervision} --- We replace lexical fuzzy matching with Sentence-BERT similarity for label generation, so extractor targets are based on semantic alignment rather than string overlap.
    
    \item \textbf{Extension 2: Multilingual Adaptation} --- We validate the framework on Italian (WITS), where richer morphology makes lexical matching particularly fragile.
\end{enumerate}

Empirically, the Italian extension is robust across configurations and shows strong data efficiency, with the 10k setup matching or outperforming larger variants.

%% ============================================================
%% SECTION 2: RELATED WORK
%% ============================================================
\section{Related Work}

\subsection{Controllable Summarization}
Traditional abstractive summarization systems~\cite{lewis2020bart, zhang2020pegasus} generate summaries without explicit control over content inclusion. Recent work has explored controllability through keywords~\cite{he2022ctrlsum}, query-focused summarization~\cite{xu2020query}, and prompt engineering~\cite{goyal2022news}.

\subsection{Salient Information Prompting}
Xu et al.~\cite{xu2024sip} introduced SigExt, a Longformer-based model trained to identify salient sentences. By injecting extracted keyphrases into LLM prompts, they achieve improved ROUGE scores while maintaining fluency. However, their reliance on lexical matching for supervision limits generalization.

\subsection{Semantic Similarity in NLP}
Sentence-BERT~\cite{reimers2019sbert} enables efficient computation of sentence embeddings for semantic similarity tasks. Multilingual variants like \texttt{paraphrase-multilingual-mpnet-base-v2} support cross-lingual applications, making them ideal for our semantic supervision approach.

\subsection{Italian Text Summarization}
Italian NLP resources are limited compared to English. The WITS dataset~\cite{casola2021wits} provides Wikipedia articles with lead sections as summaries, offering a large-scale resource for Italian summarization research.

%% ============================================================
%% SECTION 3: METHODOLOGY
%% ============================================================
\section{Methodology}

\subsection{System Architecture}

SM-SIP follows a two-stage neuro-symbolic architecture: the \textbf{Extractor (SigExt)} is a fine-tuned transformer that identifies salient content in the source document and outputs token-level binary predictions, which are aggregated to sentence-level salience scores during inference. The \textbf{Generator (LLM)} is Llama-3.1-8B-Instruct with 4-bit or 8-bit quantization~\cite{dettmers2022llm}, conditioned on extracted keyphrases via dynamic prompting to ensure controlled generation. To handle distinct linguistic challenges, we employ domain-specific backbones for the \textbf{Extractor}: 
\begin{itemize}
    \item  \textbf{English (ArXiv):} We use \path{allenai/longformer-large-4096}, a long-context encoder suitable for token-classification in technical domains;
    \item  \textbf{Italian (WITS):} We use \path{markussagen/xlm-roberta-longformer-base-4096}; although the backbone supports long-context inputs, our training setup uses a maximum sequence length of 2048 tokens.
\end{itemize}

\subsection{Extension 1: Semantic Supervision}
\label{sec:semantic}

To overcome the limitation of fuzzy matching, we introduce Semantic Supervision. We generate training labels using Sentence-BERT embedding with the multilingual model \texttt{paraphrase-multilingual-mpnet-base-v2}. A source sentence  $\vec{p}$ is assigned a positive label if its cosine similarity with any summary sentence $\vec{q}$ exceeds the similarity threshold $\theta$:

$$\text{Label}(p) = \mathbbm{1}\left( \max_{q \in S_{\text{ref}}} \cos(\vec{p}, \vec{q}) > \theta \right)$$

In the Italian extension we investigate $\theta \in \{0.60, 0.65, 0.70\}$ to identify the optimal balance between noise reduction and paraphrase retention, while in the English extension we use $\theta=0.60$. This implementation bridges the practical engineering components (model architectures and embeddings) with the theoretical foundation of semantic supervision, providing a coherent connection between system design and conceptual motivation.

\subsection{Extension 2: Italian Multilingual Adaptation} \label{sec:italian}
We validate the approach on Italian using the WITS dataset. Since traditional fuzzy matching fails with Italian morphology (e.g. verb conjugation, subject omitted, or article agreement) semantic supervision is essential.

\subsubsection{Prompt Engineering}
To enforce controllability, we designed a rigorous "Source-Aware" prompt template structured in three tiers:
\begin{itemize}
    \item \textbf{Security Rules:} Frame the task as objective summarization to bypass safety filters on sensitive topics.
    \item \textbf{Accuracy Rules:} Explicitly forbid external knowledge to mitigate hallucinations ("If the source doesn't mention it, omit it").
    \item \textbf{Writing Rules:} Impose stylistic constraints (e.g., max 120 words, no lists) to prevent the model from merely copying source segments.
\end{itemize}
This structured prompting was critical to achieving a high Faithfulness score by grounding the LLM strictly in the extracted keyphrases. However, the presence of the Security Rules may inadvertently provoke the model to refuse the task when using a Greedy Configuration.

\subsection{Experimental Setup} \label{sec:exp_config}
\textbf{1) Training Configuration} 
We train SigExt variants across two scales to test data efficiency:
\begin{itemize}
    \item \textbf{English (ArXiv):} 1k and 5k samples ($\theta=0.60$) to analyze semantic learning on limited data.
    \item \textbf{Italian (WITS):} 10k and 25k samples ($\theta \in \{0.60, 0.65, 0.70\}$).
\end{itemize}
 The English training uses per-device batch size 1, gradient accumulation 16, learning rate $1 \times 10^{-5}$, BF16 precision, and positive class weight 10.0. The Italian training uses per-device batch size 2, gradient accumulation 4, learning rate $2 \times 10^{-5}$, FP16 precision, and positive class weight 10.0. For \textit{Extension II}, we compare \textbf{Greedy Decoding} against \textbf{Stochastic Sampling} (Temperature = $0.1$, $0.2$, Top\_p = $0.9$), applying Repetition Penalty of $1.10$ to prevent cyclic loops.\\
 \textbf{2) Inference Strategies:} We assess the model using two prompting strategies to test in-context learning capabilities:
\begin{itemize}
    \item \textbf{Zero-Shot:} Direct prompting without examples, relying solely on the extracted keyphrases.
    \item \textbf{Few-Shot:} Providing 3 demonstration examples to guide stylistic alignment.
\end{itemize}
 \textbf{3) Evaluation Metrics} Beyond standard metrics such as ROUGE and BERTScore, we introduce additional metrics to ensure semantic steerability: 
 \begin{itemize}
     \item \textbf{Key Information Retention (KIR):} The percentage of extracted keyphrases successfully incorporated into the summary;
    \item  \textbf{Abstraction Score}: Computed as $(1 - \text{n-gram overlap})$, it measures whether the model synthesizes concepts rather than performing verbatim extraction;
    \item \textbf{LLM-as-a-Judge:} We employ a Qwen2.5-14B judge~\cite{goyal2022news} within a G-EVAL framework~\cite{microsofteval} and Chain-of-Thought reasoning~\cite{CoT} to score Faithfulness, Completeness, Conciseness, and Abstraction on a 1-5 scale.
 \end{itemize}
\begin{table}[h]
\centering
\caption{Training Configurations for Extension I (English) and II (Italian)}
\label{tab:combined_configs}
\resizebox{\columnwidth}{!}{%
\begin{tabular}{l c c l}
\toprule
\textbf{Config} & \textbf{Samples} & \textbf{$\theta$} & \textbf{HuggingFace ID} \\
\midrule
\multicolumn{4}{l}{\textit{\textbf{Extension I: English (ArXiv)}}} \\
1k-60t & 1,000 & 0.60 & LookUpMark/sigext-arxiv-en-1k-060t \\
5k-60t & 5,000 & 0.60 & LookUpMark/sigext-arxiv-en-5k-060t \\
\midrule
\multicolumn{4}{l}{\textit{\textbf{Extension II: Italian (WITS)}}} \\
10k-60t & 10,000 & 0.60 & LookUpMark/sigext-wits-it-10k-060t \\
25k-60t & 25,000 & 0.60 & LookUpMark/sigext-wits-it-25k-060t \\
25k-65t & 25,000 & 0.65 & LookUpMark/sigext-wits-it-25k-065t \\
25k-70t & 25,000 & 0.70 & LookUpMark/sigext-wits-it-25k-070t \\
\bottomrule
\end{tabular}%
}
\end{table}


%% ============================================================
%% SECTION 4: RESULTS
%% ============================================================
\section{Experimental Results}
\subsection{Quantitative Performance}
We evaluate SM-SIP across both extensions using the methodology described in Section \ref{sec:exp_config}. Table \ref{tab:results_unified_full} summarizes the key performance metrics.\\
\textbf{1) English ArXiv:} The English model exhibits a strong scaling trend. Increasing the training data from $1$k to $5$k yields a performance boost, with BERTScore rising to \textbf{0.886} and KIR from $73.10\%$ to $89.40\%$. Furthermore, \textbf{Few-Shot} prompting outperforms \textbf{Zero-Shot} prompting, suggesting that, in technical domains, demonstrations are crucial for alignment between the model and the expected output.\\
\textbf{2) Italian WITS:} Conversely, the Italian model reveals a \textbf{Data Efficiency Paradox}. The smallest configuration (\textbf{10k samples}) matches or outperforms the larger ones ($25$k samples), likely due to cleaner semantic labels, making the signature extraction task learnable with fewer examples. Here, the prompting trade-off is inverted: \textbf{Zero-Shot} achieves the highest KIR (49.1\%), while Few-Shot tends to prioritize fluency over strict extraction.
\begin{table}[h]
\centering
\caption{Unified Results: English (ArXiv) vs Italian (WITS)}
\label{tab:results_unified_full}
\resizebox{\columnwidth}{!}{%
\begin{tabular}{|c|c|c|c|c|c|c|c|}
\hline
\multirow{2}{*}{\textbf{Config}} & \multirow{2}{*}{\textbf{Q-bit}} & \multicolumn{3}{c|}{\textbf{Zero-Shot}} & \multicolumn{3}{c|}{\textbf{Few-Shot}} \\
\cline{3-8} 
 &  & \textbf{BERT} & \textbf{ROUGE} & \textbf{KIR (\%)} & \textbf{BERT} & \textbf{ROUGE} & \textbf{KIR (\%)} \\
\hline
\multicolumn{8}{|c|}{\textit{\textbf{Extension 1: English (ArXiv)}}} \\
\hline
\multirow{2}{*}{1k-60t} & 4 & 0.8654 & 0.3842 & 68.20 & 0.8712 & 0.4015 & 72.45 \\
 & 8 & 0.8659 & 0.3850 & 69.12 & 0.8724 & 0.4033 & 73.10 \\
\hline
\multirow{2}{*}{5k-60t} & 4 & 0.8788 & 0.4122 & 84.50 & 0.8845 & 0.4290 & 88.15 \\
 & 8 & 0.8795 & 0.4135 & 85.22 & \textbf{0.8862} & \textbf{0.4312} & \textbf{89.40} \\
\hline
\multicolumn{8}{|c|}{\textit{\textbf{Extension 2: Italian (WITS)}}} \\
\hline
\multirow{2}{*}{10k-60t} & 4 & 0.6611 & 0.2109 & 48.96 & 0.6619 & 0.2104 & 45.57 \\
 & 8 & 0.6592 & 0.2105 & \textbf{49.08} & \textbf{0.6649} & \textbf{0.2176} & 43.96 \\
\hline
\multirow{2}{*}{25k-60t} & 4 & 0.6609 & 0.2093 & 47.12 & 0.6628 & 0.2133 & 43.78 \\
 & 8 & 0.6610 & 0.2120 & 45.86 & 0.6609 & 0.2087 & 43.43 \\
\hline
\end{tabular}%
}
\end{table}

\subsection{Qualitative Analysis}
To assess quality beyond overlap-based metrics, we employ LLM-as-Judge (Qwen2.5-14B-Instruct-AWQ) with a Source-Aware prompt structured with the G-EVAL framework and Chain-of-Thought. The Italian model achieves high scores across multiple dimensions (Faithfulness: 4.17/5, Completeness: 4.66/5). However, the notably lower Abstraction Score (1.47/5) suggests that the model tends to reformulate source content with minor lexical changes rather than performing true semantic synthesis—a known limitation of current LLMs, where progress in fluency often comes at the expense of factual consistency~\cite{ladhak2022faithful, tang2024faithfulness}.

\subsection{Impact of decoding strategies}
Table \ref{tab:results_summary} summarizes the impact of decoding strategies on summary quality. While the \textbf{Greedy Decoding} strategy maximizes strict adherence (KIR 56.5\%), \textbf{Stochastic Sampling} with low temperature ($T=0.1$ and $T=0.2$) obtains higher overall scores. 

\begin{table}[h]
\caption{Comparative performance metrics across different decoding configurations (J denotes Judge scores).}
\label{tab:results_summary}
\centering
\small
\begin{tabularx}{\columnwidth}{lXXXXcc}
\toprule
\textbf{Config.} & \textbf{Faith. (J)} & \textbf{Compl. (J)} & \textbf{Conc. (J)} & \textbf{Abst. (J)} & \textbf{BERT} & \textbf{KIR\%} \\
\midrule
Config. 1 & 4.167 & 4.656 & 3.188 & 1.469 & 0.662 & \textbf{56.5} \\
Config. 2 & \textbf{4.337} & \textbf{4.667} & \textbf{3.323} & \textbf{2.042} & 0.662 & 51.2 \\
Config. 3 & 4.306 & \textbf{4.667} & 3.176 & 1.776 & 0.662 & 53.8 \\
\bottomrule
\end{tabularx}
\end{table}

%% ============================================================
%% SECTION 5: DISCUSSION
%% ============================================================
\section{Discussion}
The results suggest compact but relevant lessons for controllable multilingual summarization.

\subsection{The Paradox of Data Efficiency}
The most notable result is that the 10k configuration is as good as or better than larger variants. A plausible explanation is that semantic supervision already yields sufficiently clean targets, so the extractor converges quickly and gains from additional data are limited. Another possibility is that larger sets introduce more borderline labels, increasing supervision noise. Practically, this supports a low-cost training strategy: smaller, cleaner datasets may be enough to obtain strong controllability.

\subsection{Generator Dominance}
Performance differences remain narrow (BERTScore spread \(\approx 0.6\%\)), indicating that the generator largely determines fluency and overall quality. However, KIR confirms that SigExt is still essential: its value is not stylistic improvement, but coverage control. In other words, the extractor steers \textit{what} is preserved, while the LLM determines \textit{how} it is written.

\subsection{Robustness and Deployment Considerations}
Near-equivalent behavior between 4-bit and 8-bit settings is favorable for deployment, since memory can be reduced with minimal quality impact. Robustness across threshold values ($\theta \in \{0.60, 0.65, 0.70\}$) also reduces hyperparameter sensitivity. Together, these factors make the pipeline easier to move from experimentation to production-like settings.

\subsection{Zero-Shot vs Few-Shot Trade-offs}
Zero-shot and few-shot optimize different goals. Zero-shot improves KIR (better adherence to extracted content), while few-shot tends to improve style-oriented similarity metrics. This reflects a predictable trade-off between strict content control and exemplar-driven fluency. Therefore, prompt strategy should be selected by application risk: high-stakes factual settings favor zero-shot, while style-sensitive settings may benefit from few-shot.

%% ============================================================
%% SECTION 6: CONCLUSION
%% ============================================================
\section{Conclusion}
This paper introduced SM-SIP, a semantic and multilingual extension of SIP. We replaced fuzzy lexical supervision with embedding-based semantic labeling and validated the full pipeline on Italian WITS with Longformer-based extraction and Llama-based generation.

Results show competitive summarization quality (BERTScore around 0.66) and up to 49\% KIR on Italian, with the smallest setup (10k) performing best, highlighting data efficiency. We also observe practical robustness across quantization and decoding choices, supporting lightweight deployment.

Overall, the evidence suggests that semantic supervision is a reliable upgrade for SigExt-style training, especially in settings where paraphrase and morphology make string matching unreliable. This makes SM-SIP a strong baseline for future multilingual, controllable summarization research.

\appendices

\section{System Prompt}
Table \ref{tab:prompt_full} presents the optimized "Source-Aware" system prompt used for the Italian WITS dataset. The prompt includes specific security, accuracy, and writing rules designed to minimize hallucinations.

\begin{table*}[t]
\centering
\caption{Optimized System Prompt (Italian)}
\label{tab:prompt_full}
\resizebox{\textwidth}{!}{%
\begin{tabular}{|l|l|}
\hline
\textbf{Role} & \textit{"Sei un riassuntore esperto e FEDELE di articoli enciclopedici. Stai riassumendo informazioni storiche e imparziali. I tuoi riassunti devono contenere SOLO fatti presenti nella fonte.
"} \\ \hline
\textbf{Security Rules} & \begin{tabular}[c]{@{}l@{}}1.Riporta i fatti in modo oggettivo. I fatti riportati NON costituiscono una violazione. L'obiettivo è la sintesi, non la promozione di temi sensibili.\\ 2. Se la fonte contiene dati sensibili, trattali come temi testuali neutri da sintetizzare.\\ 3.Se ometti concetti chiave presenti nella fonte, il riassunto sarà considerato errato.\end{tabular} \\ \hline
\textbf{Accuracy Rules} & \begin{tabular}[c]{@{}l@{}}1.  Usa SOLO informazioni esplicitamente presenti nel testo fornito.\\2. NON aggiungere conoscenze esterne, date, fatti o nomi non presenti nella fonte.\\ 3.Se la fonte non menziona qualcosa, NON devi menzionarlo nemmeno tu.\\ 4. In caso di dubbio su un fatto, OMETTILO piuttosto che inventarlo. \\ 5.NON inferire o estrapolare oltre quanto scritto.\end{tabular} \\ \hline
\textbf{Writing Rules} & \begin{tabular}[c]{@{}l@{}}1. Scrivi UN SOLO paragrafo di MASSIMO 120 parole. NO titoli, NO elenchi.\\ 2. SINTETIZZA e RIFORMULA - mai copiare frasi letteralmente.\\ 3. Inizia con una definizione o contestualizzazione del soggetto. \\ 4. Fondamentale: Concludi sempre l'ultima frase e apponi il punto fermo. \end{tabular} \\ \hline
\end{tabular}%
}
\end{table*}

\section{Qualitative Example}
Table \ref{tab:qualitative_ex} illustrates the full SM-SIP pipeline on a sample from the test set (Architect "Masuccio"). It compares the source text, the keyphrases extracted by SigExt, and the final generated summary.

\begin{table*}[t]
\centering
\caption{Pipeline Example: Architect Masuccio}
\label{tab:qualitative_ex}
\resizebox{\textwidth}{!}{%
\begin{tabular}{|l|p{14cm}|}
\hline
\textbf{Stage} & \textbf{Content (Translated/Original)} \\ \hline
\textbf{Source (Excerpt)} & \textit{"Mancano notizie sulla prima parte della vita e sulla sua formazione iniziale. Entrato nell'ordine dei gesuiti venne inviato dall'ordine a Roma nel 1597, per perfezionare i suoi studi di architettura [...]"} \\ \hline
\textbf{SigExt Extraction} & 1. \textit{"Entrato nell'ordine dei gesuiti venne inviato [...] a Roma nel 1597"}\\
 & 2. \textit{"per perfezionare i suoi studi di architettura"} \\
 & 3. \textit{"Altrettanto evidente l'influsso dello stile michelangiolesco"} \\ \hline
\textbf{Generated Summary} & \textit{"Masuccio entrò nell'ordine dei gesuiti e venne inviato a Roma nel 1597 per perfezionare i suoi studi di architettura. Probabilmente ebbe modo di avviare la sua formazione in un ambiente culturale vivace [...] Nel 1616 venne espulso dall'ordine e divenne architetto del Senato messinese..."} \\ \hline
\textbf{Metrics} & \textbf{Faithfulness:} 5/5 \quad \textbf{Completeness:} 5/5 \quad  \textbf{Conciseness:} 4/5 \quad \textbf{Abstraction:} 1/5 \quad \textbf{BERT Score:} 0.648 \quad \textbf{KIR:} 25.9\% \\ \hline
\end{tabular}%
}
\end{table*}

\section*{Acknowledgment}

This work was conducted as part of the Deep Natural Language Processing course at the Politecnico di Torino.

\begin{thebibliography}{00}
\bibitem{xu2024sip} L. Xu, M. A. Karim, S. Dingliwal, and A. Elangovan, ``Salient Information Prompting to Steer Content in Prompt-based Abstractive Summarization,'' in \textit{Proc. EMNLP Industry Track}, pp. 35--49, 2024.
\bibitem{lewis2020bart} M. Lewis, Y. Liu, N. Goyal, M. Ghazvininejad, A. Mohamed, O. Levy, V. Stoyanov, and L. Zettlemoyer, ``BART: Denoising Sequence-to-Sequence Pre-training for Natural Language Generation, Translation, and Comprehension,'' in \textit{Proc. ACL}, 2020.
\bibitem{zhang2020pegasus} J. Zhang, Y. Zhao, M. Saleh, and P. J. Liu, ``PEGASUS: Pre-training with Extracted Gap-sentences for Abstractive Summarization,'' in \textit{Proc. ICML}, 2020.
\bibitem{reimers2019sbert} N. Reimers and I. Gurevych, ``Sentence-BERT: Sentence Embeddings using Siamese BERT-Networks,'' in \textit{Proc. EMNLP-IJCNLP}, 2019.
\bibitem{casola2021wits} S. Casola and A. Lavelli, ``WITS: Wikipedia for Italian Text Summarization,'' in \textit{Proc. CLiC-it}, 2021.
\bibitem{beltagy2020longformer} I. Beltagy, M. E. Peters, and A. Cohan, ``Longformer: The Long-Document Transformer,'' \textit{arXiv preprint arXiv:2004.05150}, 2020.
\bibitem{he2022ctrlsum} J. He, W. Kryscinski, B. McCann, N. Rajani, and C. Xiong, ``CTRLsum: Towards Generic Controllable Text Summarization,'' in \textit{Proc. EMNLP}, 2022.
\bibitem{xu2020query} Y. Xu and M. Lapata, ``Coarse-to-Fine Query Focused Multi-Document Summarization,'' in \textit{Proc. EMNLP}, pp. 3632--3645, 2020.
\bibitem{goyal2022news} T. Goyal, J. J. Li, and G. Durrett, ``News Summarization and Evaluation in the Era of GPT-3,'' \textit{arXiv preprint arXiv:2209.12356}, 2022.
\bibitem{dettmers2022llm} T. Dettmers, M. Lewis, Y. Belkada, and L. Zettlemoyer, ``LLM.int8(): 8-bit Matrix Multiplication for Transformers at Scale,'' in \textit{Proc. NeurIPS}, 2022.
\bibitem{maynez2020faithfulness} J. Maynez, S. Narayan, B. Bohnet, and R. McDonald, ``On Faithfulness and Factuality in Abstractive Summarization,'' in \textit{Proc. ACL}, pp. 1906--1919, 2020.
\bibitem{cao2022hallucination} M. Cao, Y. Dong, and J. Cheung, ``Hallucinated but Factual! Inspecting the Factuality of Hallucinations in Abstractive Summarization,'' in \textit{Proc. ACL}, pp. 3340--3354, 2022.
\bibitem{ladhak2022faithful} F. Ladhak, E. Durmus, H. He, C. Cardie, and K. McKeown, ``Faithful or Extractive? On Mitigating the Faithfulness-Abstractiveness Trade-off in Abstractive Summarization,'' in \textit{Proc. ACL}, pp. 1410--1421, 2022.
\bibitem{tang2024faithfulness} L. Tang, S. Goyal, A. Fabbri, P. Laban, J. Xu, Y. Dong, K. McKeown, and W. Yin, ``Understanding Factual Errors in Summarization: Errors, Summarizers, Datasets, Error Detectors,'' in \textit{Proc. ACL}, 2024.
\bibitem{microsofteval} Y. Liu, D. Iter, Y. Xu, S. Wang, R. Xu, and C. Zhu, ``G-EVAL: NLG Evaluation using GPT-4 with Better Human Alignment,'' in \textit{Proc. EMNLP}, 2023.
\bibitem{CoT} S. Saha, M. Ghazvininejad, J. Weston, and T. Wang, ``Learning to Plan \& Reason for Evaluation with Thinking-LLM-as-a-Judge,'' \textit{arXiv preprint arXiv:2501.18099}, 2025.
\end{thebibliography}

\end{document}
